% Method Section
% This file contains the detailed method description

\section{Method}
\label{sec:method}

% In this section, we first review the formulation of Rotary Position Embedding (RoPE) and its conventional extension to 2D images. We then present our Spiral RoPE and analyze its properties.

\subsection{Preliminaries: Rotary Position Embedding}

\subsubsection{1D RoPE}

Rotary Position Embedding \cite{suRoFormerEnhancedTransformer2023} encodes positional information through rotation operations. Given a $d$-dimensional query or key vector $\mathbf{x} \in \mathbb{R}^d$ at position $m$, RoPE applies a block-diagonal rotation matrix: 
\begin{equation}
    \text{RoPE}(\mathbf{x}, m) = \mathbf{R}_m \mathbf{x}
\end{equation}
where $\mathbf{R}_m$ is a block-diagonal matrix composed of $d/2$ 2D rotation matrices:
\begin{equation}
    \mathbf{R}_m = \begin{pmatrix}
        \mathbf{R}_m^{(0)} & & \\
        & \ddots & \\
        & & \mathbf{R}_m^{(d/2-1)}
    \end{pmatrix}
\end{equation}
with each $2 \times 2$ block being:
\begin{equation}
    \mathbf{R}_m^{(t)} = \begin{pmatrix}
        \cos(m\theta_t) & -\sin(m\theta_t) \\
        \sin(m\theta_t) & \cos(m\theta_t)
    \end{pmatrix}
\end{equation}

The frequency $\theta_t$ is typically set as $\theta_t = \theta_{\text{base}}^{-t/(d/2)}$ for $t = 0, 1, \ldots, d/2-1$, where $\theta_{\text{base}}$ is the base frequency that is commonly set to 10,000. This geometric progression of frequencies allows RoPE to encode positional dependencies at multiple scales.

A key property of RoPE is that the dot product between rotated vectors depends only on their relative position:
\begin{equation}
    \langle \text{RoPE}(\mathbf{q}, m), \text{RoPE}(\mathbf{k}, n) \rangle = \mathbf{q}^\top \mathbf{R}_{m-n} \mathbf{k}
    \label{eq:rope_relative}
\end{equation}
This property enables RoPE to implicitly encode relative positional information when calculating the attention matrix in the self-attention mechanism.

\subsubsection{Axial 2D RoPE}

RoPE was originally designed for 1D sequences. For 2D images, the standard approach is to extend RoPE by decomposing the spatial position into x and y coordinates. Given a patch at position $(p_x, p_y)$, the $d$-dimensional embedding $\mathbf{x} \in \mathbb{R}^d$ is split into two halves:
\begin{align}
    \mathbf{x}^{\text{(x)}} &= (x_0, x_1, \ldots, x_{d/2-1}) \\
    \mathbf{x}^{\text{(y)}} &= (x_{d/2}, x_{d/2+1}, \ldots, x_{d-1})
\end{align}

The first half is rotated using the x-coordinate, and the second half using the y-coordinate:
\begin{equation}
    \text{RoPE-2D}(\mathbf{x}, p_x, p_y) = \begin{pmatrix}
        \mathbf{R}_{p_x} \mathbf{x}^{\text{(x)}} \\
        \mathbf{R}_{p_y} \mathbf{x}^{\text{(y)}}
    \end{pmatrix}
\end{equation}

The rotation matrices $\mathbf{R}_{p_x}$ and $\mathbf{R}_{p_y}$ are typically constructed using the same set of frequencies. Each sub-embedding
$\mathbf{x}^{\text{(x)}}$ and $\mathbf{x}^{\text{(y)}}$ has dimensionality $d/2$ and thus can only accommodate $d/4$ distinct rotary frequencies, \emph{\emph{i.e.}}, half of those available in the original 1D RoPE.

Despite the good simplicity of transferring 1D RoPE to 2D inputs, this axial design only encodes positional relationships along the axis-aligned directions, leading to an inherent limitation. The x-coordinate and y-coordinate are processed independently, without any explicit encoding of diagonal or other directional relationships. Consider patches at positions $\mathbf{p}_1 = (0, 0)$ and $\mathbf{p}_2 = (1, 1)$, which are diagonally adjacent. With conventional 2D RoPE, the relative position encoding in the x-dimensions reflects $\Delta x = 1$, and in the y-dimensions reflects $\Delta y = 1$, but there is no direct encoding of the diagonal relationship.

\subsection{Spiral RoPE}

To address this limitation, we propose to extend 2D RoPE by applying additional \textbf{\textit{spatial rotations}} on the 2D panel. For example, we can directly encode the diagonal displacement $\sqrt{2}$ along the 45-degree direction, which can provide richer positional information. And more generally, we can encode positions along multiple uniformly distributed directions, which provides the model with a more complete description of spatial relationships.

Formally, let $K$ be the number of directions we want to encode. We uniformly distribute $K$ directions in the angular range $[0, \pi)$:
\begin{equation}
    \phi_k = \frac{k \pi}{K}, \quad k = 0, 1, \ldots, K-1
\end{equation}

Note that we only need to cover $[0, \pi)$ rather than $[0, 2\pi)$ because directions $\phi$ and $\phi + \pi$ are parallel, and RoPE naturally handles both positive and negative positions through its rotation mechanism.

For each direction $\phi_k$, we define a unit direction vector:
\begin{equation}
    \mathbf{u}_k = (\cos\phi_k, \sin\phi_k)
\end{equation}

Given a patch at 2D position $\mathbf{p} = (p_x, p_y)$, we compute its projected position along direction $k$ as:
\begin{equation}
    t_k(\mathbf{p}) = \mathbf{p} \cdot \mathbf{u}_k = p_x \cos\phi_k + p_y \sin\phi_k
    \label{eq:projection}
\end{equation}

We partition the $d$-dimensional embedding into $K$ equal groups, each with $d/K$ dimensions. The $k$-th group, denoted $\mathbf{x}^{(k)}$, is rotated using the projected position $t_k$:
\begin{equation}
    \text{Spiral RoPE}(\mathbf{x}, \mathbf{p}) = \begin{pmatrix}
        \mathbf{R}_{t_0(\mathbf{p})}^{(0)} \mathbf{x}^{(0)} \\
        \mathbf{R}_{t_1(\mathbf{p})}^{(1)} \mathbf{x}^{(1)} \\
        \vdots \\
        \mathbf{R}_{t_{K-1}(\mathbf{p})}^{(K-1)} \mathbf{x}^{(K-1)}
    \end{pmatrix}
\end{equation}
where $\mathbf{R}_{t}^{(k)}$ denotes the rotation matrix for direction $k$ with its specific frequency assignment.


A critical design choice in Spiral RoPE is how to assign rotation frequencies $\theta_i$ to different directions. A naive approach would independently assign $d/(2K)$ frequencies to each direction. However, this limits the number of distinct frequencies as $K$ increases, reducing the model's ability to encode positions at multiple scales.

To maximize frequency diversity, we adopt an interleaved frequency assignment strategy. We first compute a pool of $d/4$ base frequencies using the same formula as axial 2D RoPE:
\begin{equation}
    \theta_t = \theta_{\text{base}}^{-t/(d/4)}, \quad t = 0, 1, \ldots, d/4-1
\end{equation}
This gives us the same number of distinct frequencies as the conventional 2D RoPE baseline. 

To distribute these frequencies across directions, we adopt a grouped interleaved assignment strategy. We first group adjacent frequencies into pairs: $(\theta_0, \theta_1), (\theta_2, \theta_3), \ldots, (\theta_{d/4-2}, \theta_{d/4-1})$. We then assign these pairs to direction groups in a round-robin fashion. Specifically, we pair directions that are 90° apart (\emph{i.e.}, perpendicular), since they encode orthogonal spatial relationships. For direction pair $(k, k + K/2)$ where $k < K/2$, we assign frequency pairs:
\begin{equation}
    \Theta^{(k)} = \Theta^{(k+K/2)} = \{\theta_{2k}, \theta_{2k+1}, \theta_{2k+K}, \theta_{2k+K+1}, \ldots\}
\end{equation}

For example, with $K=4$ directions (0$^\circ$, 45$^\circ$, 90$^\circ$, 135$^\circ$) and $d=32$, the frequency assignment is:
\begin{itemize}
    \item $\theta_0, \theta_1, \theta_4, \theta_5$ are assigned to directions 0$^\circ$ and 90$^\circ$
    \item $\theta_2, \theta_3, \theta_6, \theta_7$ are assigned to directions 45$^\circ$ and 135$^\circ$
\end{itemize}

\cref{fig:teaser} visualizes the frequency assignment by plotting each frequency $\pm\theta_t$ as a point in the 2D frequency plane along the direction it encodes. For traditional 2D axial RoPE, all frequencies lie on the horizontal ($0^\circ$) and vertical ($90^\circ$) axes, corresponding to frequency vectors $(\pm\theta_t, 0)$ and $(0, \pm\theta_t)$. In contrast, Spiral RoPE with $K=8$ distributes frequencies across $K/2=4$ pairs of perpendicular directions: $0^\circ/90^\circ$, $22.5^\circ/112.5^\circ$, $45^\circ/135^\circ$, and $67.5^\circ/157.5^\circ$, which forms a \textit{spiral} pattern. This visualization illustrates how Spiral RoPE expands directional coverage beyond the coordinate axes while maintaining the same total number of frequencies.

This grouped interleaved assignment ensures that: (1) we utilize all $d/4$ distinct frequencies, matching the capacity of axial 2D RoPE; (2) perpendicular direction pairs share the same frequency set, analogous to how axial 2D RoPE uses the same frequencies for both x and y axes; and (3) each direction receives a mix of low and high frequencies from adjacent pairs, enabling multi-scale positional encoding within each direction.

Spiral RoPE preserves the relative position encoding property of RoPE: for patches at positions $\mathbf{p}$ and $\mathbf{p}'$, since $t_k(\mathbf{p}) - t_k(\mathbf{p}') = t_k(\mathbf{p} - \mathbf{p}')$, the attention score contribution depends only on the relative position $\mathbf{p} - \mathbf{p}'$. Furthermore, by uniformly distributing directional encodings, Spiral RoPE provides more isotropic spatial awareness compared to axial 2D RoPE, which treats horizontal and vertical directions preferentially. Note that Spiral RoPE can be implemented efficiently by precomputing the rotation matrices for all positions and directions, with minimal computational overhead compared to axial 2D RoPE.

To further illustrate the representation power of Spiral RoPE, we perform a 2D Fourier analysis comparing axial 2D RoPE with Spiral RoPE. We consider $d=1024$, which yields $d/4=256$ distinct RoPE frequencies. Given a binary image, we first apply 2D FFT to obtain its frequency-domain representation. We then retain only the frequency components that can be represented by each RoPE variant: for axial 2D RoPE, we keep frequencies along the horizontal and vertical axes $(\pm\theta_t, 0)$ and $(0, \pm\theta_t)$; for Spiral RoPE, we keep frequencies along all $K=8$ directions as illustrated in \cref{fig:teaser}. All other frequency components are masked to zero. Finally, we apply 2D iFFT to reconstruct the image and compare the result with the original input. \cref{fig:fourier} shows the reconstruction results for two input images: a single point and a circle on the $64 \times 64$ grid. The axial 2D RoPE reconstruction exhibits visible artifacts along the horizontal and vertical axes, as it can only represent frequency components in these two directions. In contrast, Spiral RoPE produces reconstructions that are closer to the original images, demonstrating its ability to capture a broader range of spatial patterns through multi-directional frequency coverage.

